\section{Disjoint support components: a dependent consequence and scope check}
\label{bs:section}

The aggregate-core estimates suggest covering prime-disjoint source
components separately. That argument is valid, but for components with
$o(\log M)$ primes it does not produce a previously untreated integer
class: an elementary smooth-number count and the inherited quadratic
theorem already give an eventual exact $2n$ bound. We record the stronger
dependent consequence to make this limitation explicit.

\begin{lemma}[A uniform count on a small prime support]
\label{bs:rankin}
For every finite prime set $Q$ with $|Q|\le k$, the number $\Psi_Q(M)$
of $Q$-smooth positive integers at most $M$ satisfies
\begin{equation}
 \Psi_Q(M)\le M^{1/4}\exp(Ck^{3/4}),
 \qquad C=\frac4{3(1-2^{-1/4})}.
 \label{bs:smooth-count}
\end{equation}
The bound is uniform in the actual prime values.
\end{lemma}

\begin{proof}
Rankin's inequality at exponent $1/4$ gives
\[
 \Psi_Q(M)\le
 M^{1/4}\sum_{\substack{h\ge1\\h\ Q\text{-smooth}}}h^{-1/4}
 =M^{1/4}\prod_{p\in Q}(1-p^{-1/4})^{-1}.
\]
If the primes are listed as $p_1<\cdots<p_f$, then $p_j\ge j+1$.
Using $-\log(1-x)\le x/(1-x)$, we obtain
\[
 \begin{split}
 \log\prod_{p\in Q}(1-p^{-1/4})^{-1}
 &\le\frac1{1-2^{-1/4}}\sum_{j=1}^f(j+1)^{-1/4}\\
 &\le\frac4{3(1-2^{-1/4})}k^{3/4}.
 \end{split}
\]
The empty support also satisfies the bound, since it contains only
the smooth integer one.
\end{proof}

\begin{corollary}[An inherited exact bound for disjoint small supports]
\label{bs:exact}
Let $k(M)=o(\log M)$. Suppose the prime support of a distinct source set
$A\subseteq[2,M]$ can be partitioned into disjoint blocks
$P_1,\ldots,P_s$, each of size at most $k(M)$, and every source has
\emph{all} its prime factors in a single block. Then, for all sufficiently
large $M$, uniformly in the partition, the prime values, and every
positive $M$-term interval,
\[
 \tau(A,I)\le2|A|.
 \]
\end{corollary}

\begin{proof}
Group the sources as $A_i$, supported on $P_i$, and discard empty
groups. Let $n_i=|A_i|$. Lemma~\ref{bs:rankin} gives, uniformly in $i$,
\[
 n_i\le M^{1/4}\exp(Ck(M)^{3/4})
      =M^{1/4+o(1)}<\sqrt M/4
\]
for large $M$. Hence $M\ge16n_i^2$ for every group.
The ambient quadratic theorem of \cite{Wang} supplies a cover $B_i$
of $A_i$ in the original interval $I$ with $|B_i|\le2n_i$.

The union $B=\bigcup_iB_i$ covers all of $A$.
Indeed, each demanded prime belongs to a unique block. Its entire
source demand is supplied by that block's cover, which is contained
in $B$. Possible overlaps of different $B_i$ therefore cause no loss:
there is no additional source demand at that prime from another
block. Finally,
$|B|\le\sum_i|B_i|\le2\sum_i n_i=2|A|$.
\end{proof}

For comparison, the aggregate-core bound with no outside primes gives
\[
 \mathcal L_{\rm raw}(A_i,I)
 \le(2+\delta)n_i+B_{\delta,k},\qquad
 B_{\delta,k}\le D_\delta^k,
\]
and rounding adds at most $k$ positions. Applying this only to groups
with $n_i>\sqrt M/4$, while using the exact quadratic bound for smaller
groups, would give the valid estimate
\[
 \tau(A,I)\le(2+\delta)|A|
       +\frac{4|A|}{\sqrt M}(B_{\delta,k}+k).
\]
For $k=o(\log M)$ its last term is $o(|A|)$.
However, the smooth-number argument shows that eventually every group
is already small, so the fractional route is unnecessary for this class.

The component hypothesis is stronger than a bound on the number of
distinct primes in each individual source. Different sources can connect
many prime values through overlaps. It also differs from a shared-core
geometry, where common primes occur across many source groups; in that
case the prime-disjoint union argument is unavailable.
The auxiliary weighted block-gluing theorem permits such a shared load
only when its weights are diffuse, and it does not remove the arbitrary
shared-core obstruction.
